\section{新范式：Environment + Task-Reward}

本章收束到节目里的展望：新的数据核心可能从 input-output 标注，转向 environment 和 task-reward 构造。也就是说，未来训练 Agent 的核心不只是收集答案，而是构造可交互环境、定义任务、设计可验证 reward，让模型从经验中 self-improve。

\begin{figure}[H]
\centering
\includegraphics[width=0.96\textwidth]{new-agent-paradigm.png}
\caption{Agent 新范式：数据核心从 input-output 转向 environment + task-reward。自制概念图，依据 02:06:06--02:15:20 对谈内容整理。}
\end{figure}

\begin{knowledgebox}{读图：从答案数据到经验数据}
Input-output 数据教模型“看到输入给出答案”；environment + task-reward 教模型“在环境中行动并从反馈中改进”。这就是 Agent 与普通聊天模型的数据范式差异。
\end{knowledgebox}

\subsection{Rubrics as reward 与 self-improvement}

上一节讲新数据范式，本节聚焦它的核心难题：reward 从哪里来。节目提到 rubrics as reward，即用评分标准作为奖励机制。它的意义是：很多现实任务没有简单对错，但可以有分层评分标准。Agent 能不能自我提升，取决于它能否找到或构造可验证 reward，并有效利用交互经验。

\begin{warningbox}{Self-improve 的前提}
Agent 自我提升不是让模型无限自嗨。它需要可验证任务、可靠环境、可审计轨迹和防止 reward hacking 的安全机制。
\end{warningbox}

\begin{knowledgebox}{新范式对照表}
\begin{tabularx}{\textwidth}{p{0.22\textwidth}p{0.34\textwidth}X}
\toprule
范式 & 数据形态 & 主要瓶颈 \\
\midrule
Input-output & 输入和标准答案 & 标注成本和泛化边界。 \\
Trajectory & 观察、动作、工具调用、结果 & 轨迹质量和错误恢复。 \\
Environment & 可交互任务世界 & 环境构建成本和覆盖度。 \\
Task-reward & 任务定义和可验证反馈 & reward 设计和安全。 \\
Experience reuse & 利用历史经验自我提升 & 记忆、去噪和防止自嗨。 \\
\bottomrule
\end{tabularx}
\end{knowledgebox}

\subsection{Family of Agents}

reward 和 self-improvement 讨论的是单体系统如何学习，本节转向多 Agent 组织。结尾提到 Agent 像“拓展的大脑”，背后有一个军团。这个比喻说明未来 Agent 可能不是一个单体，而是一组专门化智能体：coding、search、browser、memory、safety 等共同协作。

\begin{figure}[H]
\centering
\includegraphics[width=0.92\textwidth]{family-of-agents.png}
\caption{Family of Agents：Agent 像拓展的大脑，背后是一支军团。自制概念图，依据 02:15:20--02:16:41 对谈内容整理。}
\end{figure}

\subsection{本章小结}

Agent 的新范式是把数据、环境、奖励和经验循环合在一起。真正的系统工程，是让这些循环可规模化、可验证、可安全部署。

